\documentclass[11pt,a4paper]{article}

% --- Codificación e idioma ---
\usepackage[utf8]{inputenc}
\usepackage[T1]{fontenc}
\usepackage[spanish,es-tabla]{babel}
\usepackage{lmodern}

% --- Geometría y formato general ---
\usepackage[margin=2.5cm]{geometry}
\usepackage{microtype}
\usepackage{parskip}
\usepackage{enumitem}
\setlist{topsep=2pt,itemsep=2pt,parsep=0pt}

% --- Tablas ---
\usepackage{array}
\usepackage{booktabs}

% --- Cabeceras y numeración ---
\usepackage{fancyhdr}
\pagestyle{fancy}
\fancyhf{}
\fancyhead[L]{\small Bases de Datos II}
\fancyhead[R]{\small Trabajo Práctico N.\textsuperscript{o} 4}
\fancyfoot[C]{\thepage}
\renewcommand{\headrulewidth}{0.4pt}

% --- Color y código ---
\usepackage{xcolor}
\definecolor{codebg}{RGB}{248,248,248}
\definecolor{codeframe}{RGB}{200,200,200}
\definecolor{codekw}{RGB}{0,0,180}
\definecolor{codestr}{RGB}{163,21,21}
\definecolor{codecmt}{RGB}{0,128,0}
\definecolor{notebg}{RGB}{255,250,230}
\definecolor{noteframe}{RGB}{220,180,80}

\usepackage{listings}
\lstdefinelanguage{SQLext}[]{SQL}{
    morekeywords={NUMERIC,VARCHAR,DECIMAL,DATE,TIMESTAMP,
        BIGINT,VARCHAR2,NUMBER,SCHEMA,DATABASE,
        WITH,FROM,TO,ON,AS,AND,OR,IS,NOT,NULL,
        INSERT,UPDATE,DELETE,SELECT,PROCEDURE,FUNCTION,CALL},
    sensitive=false,
    morecomment=[l]{--},
    morestring=[b]'
}

\lstdefinestyle{javastyle}{
    language=Java,
    basicstyle=\ttfamily\footnotesize,
    keywordstyle=\color{codekw}\bfseries,
    stringstyle=\color{codestr},
    commentstyle=\color{codecmt}\itshape,
    backgroundcolor=\color{codebg},
    frame=single,
    rulecolor=\color{codeframe},
    breaklines=true,
    breakatwhitespace=true,
    columns=fullflexible,
    keepspaces=true,
    showstringspaces=false,
    captionpos=b,
    xleftmargin=0.4em,
    xrightmargin=0.4em,
    aboveskip=8pt,
    belowskip=8pt,
    tabsize=2
}

\lstdefinestyle{sqlstyle}{
    language=SQLext,
    basicstyle=\ttfamily\footnotesize,
    keywordstyle=\color{codekw}\bfseries,
    stringstyle=\color{codestr},
    commentstyle=\color{codecmt}\itshape,
    backgroundcolor=\color{codebg},
    frame=single,
    rulecolor=\color{codeframe},
    breaklines=true,
    breakatwhitespace=true,
    columns=fullflexible,
    keepspaces=true,
    showstringspaces=false,
    captionpos=b,
    xleftmargin=0.4em,
    xrightmargin=0.4em,
    aboveskip=8pt,
    belowskip=8pt,
    tabsize=2
}

\lstdefinestyle{shellstyle}{
    basicstyle=\ttfamily\footnotesize,
    backgroundcolor=\color{codebg},
    frame=single,
    rulecolor=\color{codeframe},
    breaklines=true,
    breakatwhitespace=true,
    columns=fullflexible,
    keepspaces=true,
    showstringspaces=false,
    xleftmargin=0.4em,
    xrightmargin=0.4em,
    aboveskip=8pt,
    belowskip=8pt
}

\lstdefinestyle{propsstyle}{
    basicstyle=\ttfamily\footnotesize,
    keywordstyle=\color{codekw}\bfseries,
    commentstyle=\color{codecmt}\itshape,
    backgroundcolor=\color{codebg},
    frame=single,
    rulecolor=\color{codeframe},
    breaklines=true,
    breakatwhitespace=true,
    columns=fullflexible,
    keepspaces=true,
    showstringspaces=false,
    xleftmargin=0.4em,
    xrightmargin=0.4em,
    aboveskip=8pt,
    belowskip=8pt,
    morecomment=[l]{\#}
}

\lstset{style=javastyle}

% --- Cajas de nota ---
\usepackage[most]{tcolorbox}
\newtcolorbox{nota}{colback=notebg,colframe=noteframe,arc=2pt,
    left=6pt,right=6pt,top=4pt,bottom=4pt,boxrule=0.6pt}

% --- Hipervínculos ---
\usepackage[hidelinks,bookmarks=true]{hyperref}

% --- Comandos auxiliares (mismos que en Prácticos 1, 2 y 3) ---
\newcommand{\itemcabecera}[1]{\noindent\textbf{#1}\par\medskip}
\newcommand{\bloqueConsigna}{\itemcabecera{Consigna}}
\newcommand{\bloqueIdea}{\itemcabecera{Idea}}
\newcommand{\bloqueSolucion}{\itemcabecera{Soluci\'on}}
\newcommand{\bloquePorque}{\itemcabecera{Por qu\'e funciona}}
\newcommand{\bloqueMySQL}{\itemcabecera{En MySQL 8+}}
\newcommand{\bloquePrueba}{\itemcabecera{Prueba}}

\title{Trabajo Práctico N.\textsuperscript{o} 4: JDBC\\[4pt]
       \large Bases de Datos II}
\author{Tomás Rodeghiero}
\date{2026}

\begin{document}

% =====================================================
% Carátula
% =====================================================
\begin{titlepage}
    \centering
    \vspace*{1.5cm}

    {\large \textbf{Universidad Nacional de Río Cuarto}}\\[6pt]
    {\normalsize Facultad de Ciencias Exactas, Físico-Químicas y Naturales}\\[4pt]
    {\normalsize Departamento de Computación}\\[4pt]
    {\normalsize Licenciatura en Ciencias de la Computación}\\[3cm]

    {\Large \textbf{Bases de Datos II}}\\[1.2cm]

    {\Large \textbf{Trabajo Práctico N.\textsuperscript{o} 4}}\\[6pt]
    {\large JDBC --- Conexión a bases desde Java}\\[3cm]

    \begin{tabular}{rl}
        \textbf{Alumno:} & Tomás Rodeghiero \\[4pt]
        \textbf{Año:}    & 2026 \\
    \end{tabular}

    \vfill
    {\small Resolución usando la API \texttt{java.sql} (JDBC) sobre las
    bases creadas en los Prácticos 1 y 3.}
\end{titlepage}

\tableofcontents
\newpage

% =====================================================
\section{Introducción}
% =====================================================

El Práctico 4 cierra el bloque de fundamentos del cuatrimestre moviéndose desde SQL puro hacia el \emph{cliente}: cómo se accede a las bases que armamos en prácticos anteriores desde una aplicación Java mediante la API \textbf{JDBC} (\emph{Java DataBase Connectivity}). Los tres ejercicios cubren los tres usos canónicos:

\begin{itemize}
    \item \textbf{Ejercicio 1}: \texttt{PreparedStatement} para alta y \texttt{ResultSet} para listado, sobre la base de facturación del Práctico 1.a (MySQL y PostgreSQL).
    \item \textbf{Ejercicio 2}: \texttt{CallableStatement} para invocar la función \texttt{fn\_actualizar\_cotizacion\_maxima} creada en el Práctico 3, sobre PostgreSQL (y el procedimiento equivalente en MySQL).
    \item \textbf{Ejercicio 3}: \texttt{DatabaseMetaData} para listar tablas, columnas, claves primarias, únicas y foráneas de la base del Práctico 1.c, en PostgreSQL y Oracle (más MySQL como caso adicional).
\end{itemize}

Cada ejercicio sigue siempre la misma estructura: \textbf{Consigna} (texto literal del PDF), \textbf{Idea} (lectura natural y concepto JDBC detrás), \textbf{Solución} (código Java con su \texttt{.properties}), \textbf{Por qué funciona} (lógica paso a paso), \textbf{En MySQL 8+} (snippet equivalente con driver y URL de MySQL) y \textbf{Prueba} (compilación e invocación). Toda la solución usa exclusivamente \texttt{java.sql} (sin Hibernate ni JPA): JDBC pelado, como pide la consigna.

\begin{nota}
\textbf{Convenciones de proyecto.} Cada ejercicio vive en su carpeta (\texttt{ejercicio-01}, \texttt{ejercicio-02}, \texttt{ejercicio-03}) con la misma estructura: \texttt{src/} (código), \texttt{lib/} (drivers JDBC), \texttt{config/} (\texttt{*.properties}) y un \texttt{README.md} con instrucciones. Los datos sensibles (host, usuario, clave) viven en los \texttt{.properties}, no en el código. Todos los recursos JDBC se cierran con \texttt{try-with-resources}; todas las sentencias usan \texttt{?} y \texttt{PreparedStatement} (sin concatenación de strings).
\end{nota}

% =====================================================
\section{Ejercicio 1: Alta y listado de clientes}
% =====================================================

\bloqueConsigna

Teniendo en cuenta el ejercicio 1 a) de la práctica 1) (base de datos de facturas/clientes), cree un programa java que implemente un Alta y listado de clientes. El programa puede ser hecho sin Interfaz gráfica, sólo en línea de comandos.

\textbf{Nota:} Resolver utilizando los motores que se utilizaron en la creación de las tablas.

\bloqueIdea

Es el uso más típico de JDBC: \texttt{PreparedStatement} con placeholders \texttt{?} para escribir (\texttt{INSERT}) y \texttt{ResultSet} para leer (\texttt{SELECT}). El esqueleto del programa separa \emph{value object} (\texttt{Cliente}), \emph{DAO} (\texttt{ClienteDao}), \emph{config} (\texttt{DbConfig}) y \emph{main} (\texttt{ClientesCliApp}); esto permite cambiar entre MySQL y PostgreSQL apenas cambiando el archivo \texttt{.properties}.

\bloqueSolucion

\textit{Estructura de paquetes} (\texttt{ar.edu.unrc.bd2.practico4.ej1}):

\begin{itemize}
    \item \texttt{Cliente}: value object inmutable con los cinco campos.
    \item \texttt{ClienteDao}: encapsula JDBC. Recibe la \texttt{Connection} desde afuera y expone \texttt{insertar} y \texttt{listar}.
    \item \texttt{DbConfig}: lee el \texttt{.properties}, abre la conexión y expone el nombre completo de la tabla.
    \item \texttt{ClientesCliApp}: el \texttt{main} con el menú de consola.
\end{itemize}

\textit{Configuración por motor (\texttt{.properties}):}

\begin{lstlisting}[style=propsstyle]
# config/db.postgres.properties
db.driver=org.postgresql.Driver
db.url=jdbc:postgresql://localhost:5432/postgres
db.user=postgres
db.password=root
db.table.cliente=practico1a.cliente
\end{lstlisting}

\begin{lstlisting}[style=propsstyle]
# config/db.mysql.properties
db.driver=com.mysql.cj.jdbc.Driver
db.url=jdbc:mysql://localhost:3306/practico1a_mysql?useSSL=false&serverTimezone=UTC
db.user=root
db.password=root
db.table.cliente=cliente
\end{lstlisting}

\textit{Apertura de conexión:}

\begin{lstlisting}
import java.sql.Connection;
import java.sql.DriverManager;
import java.sql.SQLException;

public Connection createConnection() throws SQLException, ClassNotFoundException {
    Class.forName(driverClass);
    return DriverManager.getConnection(url, user, password);
}
\end{lstlisting}

\textit{Acceso a datos: \texttt{ClienteDao}:}

\begin{lstlisting}
import java.sql.Connection;
import java.sql.PreparedStatement;
import java.sql.ResultSet;
import java.sql.SQLException;
import java.util.ArrayList;
import java.util.List;

public int insertar(Connection connection, Cliente cliente) throws SQLException {
    String sql = "INSERT INTO " + clienteTable
        + " (nro_cliente, apellido, nombre, direccion, telefono) VALUES (?, ?, ?, ?, ?)";

    try (PreparedStatement preparedStatement = connection.prepareStatement(sql)) {
        preparedStatement.setInt(1, cliente.getNroCliente());
        preparedStatement.setString(2, cliente.getApellido());
        preparedStatement.setString(3, cliente.getNombre());
        preparedStatement.setString(4, cliente.getDireccion());
        preparedStatement.setString(5, cliente.getTelefono());
        return preparedStatement.executeUpdate();
    }
}

public List<Cliente> listar(Connection connection) throws SQLException {
    String sql = "SELECT nro_cliente, apellido, nombre, direccion, telefono FROM "
        + clienteTable + " ORDER BY nro_cliente";

    List<Cliente> clientes = new ArrayList<>();

    try (PreparedStatement preparedStatement = connection.prepareStatement(sql);
         ResultSet resultSet = preparedStatement.executeQuery()) {

        while (resultSet.next()) {
            clientes.add(new Cliente(
                resultSet.getInt("nro_cliente"),
                resultSet.getString("apellido"),
                resultSet.getString("nombre"),
                resultSet.getString("direccion"),
                resultSet.getString("telefono")
            ));
        }
    }
    return clientes;
}
\end{lstlisting}

\textit{Menú de consola:}

\begin{lstlisting}
import java.sql.Connection;
import java.util.Scanner;

private static void ejecutarMenu(Scanner scanner, Connection connection,
                                 ClienteDao clienteDao) {
    boolean salir = false;
    while (!salir) {
        mostrarMenu();
        String opcion = scanner.nextLine().trim();
        switch (opcion) {
            case "1": altaCliente(scanner, connection, clienteDao); break;
            case "2": listarClientes(connection, clienteDao); break;
            case "0": salir = true; break;
            default:  System.out.println("Opcion invalida.");
        }
    }
}
\end{lstlisting}

\bloquePorque

\begin{itemize}
    \item \texttt{Class.forName(driverClass)} fuerza la carga de la clase del driver. Desde JDBC 4.x los drivers se autorregistran por \texttt{ServiceLoader}, pero llamarlo explícitamente es la forma robusta que muestra el teórico y evita problemas si \texttt{META-INF/services} no está empaquetado.
    \item \texttt{DriverManager.getConnection(url, user, password)} devuelve la \texttt{Connection}. La URL define el motor (\texttt{jdbc:postgresql} vs \texttt{jdbc:mysql}); el resto del código no cambia.
    \item \texttt{PreparedStatement} con \texttt{?} y \texttt{setInt}/\texttt{setString} envía sentencia y valores por separado: protege contra \emph{SQL Injection} y permite que el motor cachee el plan.
    \item \texttt{executeUpdate} devuelve cuántas filas modificó (se espera 1 para un alta); \texttt{executeQuery} devuelve un \texttt{ResultSet}, recorrido con \texttt{while (rs.next())}.
    \item Las llamadas a \texttt{rs.getInt("nro\_cliente")} usan \emph{nombre de columna} en vez de índice numérico: menos eficiente pero más legible y robusto frente a cambios en el orden del \texttt{SELECT}.
    \item Todo va dentro de \texttt{try-with-resources} para que \texttt{Connection}, \texttt{PreparedStatement} y \texttt{ResultSet} se cierren aunque salte una excepción.
\end{itemize}

\bloqueMySQL

La consigna ya pide ambos motores (MySQL y PostgreSQL), y el \texttt{.properties} de MySQL aparece en la \textit{Solución}. Los puntos a recordar específicos de MySQL 8+:

\begin{itemize}
    \item \textbf{Driver}: \texttt{com.mysql.cj.jdbc.Driver} del artefacto Maven \texttt{mysql:mysql-connector-j:8.x} (o el JAR equivalente bajado a \texttt{lib/}).
    \item \textbf{URL}: \texttt{jdbc:mysql://host:3306/baseDatos?useSSL=false\&serverTimezone=UTC}. Sin \texttt{serverTimezone}, MySQL 8 puede tirar \texttt{The server time zone value is unrecognized} al abrir la conexión.
    \item \textbf{Base vs esquema}: en MySQL la ``base'' es lo que JDBC llama \texttt{catalog}; no hay un esquema separado. Por eso el \texttt{.properties} apunta directamente a \texttt{practico1a\_mysql} y la columna \texttt{db.table.cliente} no necesita prefijo.
    \item \textbf{Identificadores}: si una columna o tabla choca con una palabra reservada de MySQL, va entre backticks \texttt{`...`} en vez de comillas dobles (PostgreSQL/Oracle usan \texttt{"..."}).
\end{itemize}

\bloquePrueba

\begin{lstlisting}[style=shellstyle]
# Desde practico-04/ejercicio-01
mkdir -p out
javac -d out src/main/java/ar/edu/unrc/bd2/practico4/ej1/*.java

# PostgreSQL
java -cp "out:lib/*" \
    ar.edu.unrc.bd2.practico4.ej1.ClientesCliApp \
    config/db.postgres.properties

# MySQL
java -cp "out:lib/*" \
    ar.edu.unrc.bd2.practico4.ej1.ClientesCliApp \
    config/db.mysql.properties
\end{lstlisting}

% =====================================================
\section{Ejercicio 2: Invocar la función de máxima cotización}
% =====================================================

\bloqueConsigna

Realizar un programa Java que invoque la función implementada en el ejercicio 2 de la práctica 3 (máxima cotización) y muestre los parámetros de salida.

\bloqueIdea

Es el caso clásico de \texttt{CallableStatement}: se invoca una función almacenada del servidor con un parámetro \texttt{IN} (\texttt{cod\_banco}) y un \texttt{OUT} (\texttt{cotizacion\_maxima}). En PostgreSQL la función se llama con el \emph{escape call} \texttt{\{ ? = call func(?) \}}; se registra el \texttt{OUT} con \texttt{registerOutParameter(1, Types.NUMERIC)} y se lee con \texttt{getBigDecimal(1)}.

\bloqueSolucion

\textit{Estructura del proyecto:}

\begin{itemize}
    \item \texttt{CotizacionMaximaApp}: el \texttt{main}; lee el \texttt{cod\_banco}, invoca el service, imprime el resultado y consulta la fila actualizada.
    \item \texttt{CotizacionMaximaService}: encapsula la invocación JDBC con dos caminos: \texttt{CallableStatement} (preferido) y \texttt{SELECT * FROM funcion(?)} como \emph{fallback}.
    \item \texttt{DbConfig}: lee driver, url, credenciales y las dos consultas (call y select).
\end{itemize}

\textit{Configuración PostgreSQL:}

\begin{lstlisting}[style=propsstyle]
# config/db.postgres.properties
db.driver=org.postgresql.Driver
db.url=jdbc:postgresql://127.0.0.1:5432/postgres
db.user=postgres
db.password=

# Invocacion por CallableStatement (parametro OUT en posicion 1).
db.function.call={ ? = call practico3_ej2.fn_actualizar_cotizacion_maxima(?) }

# Fallback por SELECT.
db.function.select=SELECT * FROM practico3_ej2.fn_actualizar_cotizacion_maxima(?)

# Verificacion posterior de la fila actualizada.
db.query.banco=SELECT cod_banco, nombre, cotizacion_maxima FROM practico3_ej2.banco WHERE cod_banco = ?
\end{lstlisting}

\textit{Invocación con \texttt{CallableStatement}:}

\begin{lstlisting}
import java.math.BigDecimal;
import java.sql.CallableStatement;
import java.sql.Connection;
import java.sql.SQLException;
import java.sql.Types;

private BigDecimal invocarConCallableStatement(int codBanco, Connection connection)
        throws SQLException {
    try (CallableStatement callableStatement =
             connection.prepareCall(dbConfig.getFunctionCall())) {
        callableStatement.registerOutParameter(1, Types.NUMERIC);
        callableStatement.setInt(2, codBanco);
        callableStatement.execute();
        return callableStatement.getBigDecimal(1);
    }
}
\end{lstlisting}

\textit{Fallback con \texttt{SELECT}:}

\begin{lstlisting}
import java.math.BigDecimal;
import java.sql.Connection;
import java.sql.PreparedStatement;
import java.sql.ResultSet;
import java.sql.SQLException;

private BigDecimal invocarConSelect(int codBanco, Connection connection)
        throws SQLException {
    try (PreparedStatement preparedStatement =
             connection.prepareStatement(dbConfig.getFunctionSelect())) {
        preparedStatement.setInt(1, codBanco);
        try (ResultSet resultSet = preparedStatement.executeQuery()) {
            if (!resultSet.next()) {
                throw new SQLException("La funcion no retorno filas de salida.");
            }
            return resultSet.getBigDecimal(1);
        }
    }
}

public BigDecimal invocarFuncion(int codBanco, Connection connection) throws SQLException {
    try {
        return invocarConCallableStatement(codBanco, connection);
    } catch (SQLException e) {
        return invocarConSelect(codBanco, connection);
    }
}
\end{lstlisting}

\textit{Verificación posterior:}

\begin{lstlisting}
public BancoCotizacion consultarBanco(int codBanco, Connection connection) throws SQLException {
    try (PreparedStatement preparedStatement =
             connection.prepareStatement(dbConfig.getBancoQuery())) {
        preparedStatement.setInt(1, codBanco);
        try (ResultSet resultSet = preparedStatement.executeQuery()) {
            if (!resultSet.next()) return null;
            return new BancoCotizacion(
                resultSet.getInt("cod_banco"),
                resultSet.getString("nombre"),
                resultSet.getBigDecimal("cotizacion_maxima")
            );
        }
    }
}
\end{lstlisting}

\bloquePorque

\begin{itemize}
    \item \texttt{prepareCall} es el método de \texttt{Connection} para crear un \texttt{CallableStatement}. La sintaxis estándar JDBC para llamar a una función con \texttt{OUT} es el \emph{escape call} \texttt{\{ ? = call func(?) \}}, que el driver traduce al SQL nativo del motor.
    \item \texttt{registerOutParameter(1, Types.NUMERIC)} le dice al driver que el primer placeholder es de salida y de tipo \texttt{NUMERIC} (el tipo SQL de \texttt{cotizacion\_maxima}). \texttt{Types} viene de \texttt{java.sql.Types}.
    \item \texttt{getBigDecimal(1)} lee el valor de salida en el placeholder 1 como \texttt{BigDecimal}, que es la correspondencia natural de \texttt{NUMERIC} en Java según la tabla de tipos JDBC.
    \item El \emph{fallback} con \texttt{SELECT * FROM funcion(?)} existe porque algunos drivers/versiones de PostgreSQL no aceptan el escape call para funciones escalares. \texttt{invocarFuncion} prueba primero el camino limpio y, ante una \texttt{SQLException}, recurre al \texttt{SELECT}.
\end{itemize}

\bloqueMySQL

La función de PostgreSQL del Práctico 3 (con parámetros \texttt{IN}/\texttt{OUT}) se traduce en MySQL a un \texttt{PROCEDURE} con \texttt{IN p\_cod\_banco}, \texttt{OUT p\_cotizacion\_maxima} (porque MySQL no soporta \texttt{OUT} en \texttt{FUNCTION}; solo en \texttt{PROCEDURE}). El \emph{escape call} se invoca con \texttt{\{ CALL sp(?, ?) \}}, registrando el \texttt{OUT} en la posición 2:

\begin{lstlisting}[style=propsstyle]
# config/db.mysql.properties
db.driver=com.mysql.cj.jdbc.Driver
db.url=jdbc:mysql://localhost:3306/practico3_ej2_mysql?useSSL=false&serverTimezone=UTC
db.user=root
db.password=root

# IN en posicion 1, OUT en posicion 2 (orden de la declaracion del PROCEDURE).
db.function.call={ CALL sp_actualizar_cotizacion_maxima(?, ?) }

# Verificacion
db.query.banco=SELECT cod_banco, nombre, cotizacion_maxima FROM banco WHERE cod_banco = ?
\end{lstlisting}

\begin{lstlisting}
import java.math.BigDecimal;
import java.sql.CallableStatement;
import java.sql.Connection;
import java.sql.SQLException;
import java.sql.Types;

private BigDecimal invocarConCallableStatementMySQL(int codBanco, Connection connection)
        throws SQLException {
    try (CallableStatement callableStatement =
             connection.prepareCall("{ CALL sp_actualizar_cotizacion_maxima(?, ?) }")) {
        // IN: cod_banco
        callableStatement.setInt(1, codBanco);
        // OUT: cotizacion_maxima (DECIMAL en MySQL)
        callableStatement.registerOutParameter(2, Types.DECIMAL);
        callableStatement.execute();
        return callableStatement.getBigDecimal(2);
    }
}
\end{lstlisting}

\textit{Qué cambia respecto de PostgreSQL:}

\begin{itemize}
    \item \texttt{FUNCTION} con \texttt{OUT} se transforma en \texttt{PROCEDURE}: el escape call cambia de \texttt{\{ ? = call f(?) \}} (con \texttt{OUT} en posición 1, \texttt{IN} en 2) a \texttt{\{ CALL p(?, ?) \}} (\texttt{IN} en 1, \texttt{OUT} en 2, según el orden de declaración del procedimiento).
    \item El tipo registrado es \texttt{Types.DECIMAL} (en MySQL la columna se declaró \texttt{DECIMAL(12,4)}); en PostgreSQL es \texttt{Types.NUMERIC}. Ambos se leen con \texttt{getBigDecimal}.
    \item El \emph{fallback} con \texttt{SELECT * FROM funcion(?)} \emph{no aplica} en MySQL: no hay funciones tabulares; si el \texttt{CALL} no es soportado, no hay un atajo equivalente.
\end{itemize}

\bloquePrueba

\textit{Salida esperada (PostgreSQL):}

\begin{lstlisting}[style=shellstyle]
Funcion invocada: practico3_ej2.fn_actualizar_cotizacion_maxima
Parametro IN  p_cod_banco         = 1
Parametro OUT p_cotizacion_maxima = 1245.7500
Banco actualizado -> cod_banco=1, nombre=Banco Rio, cotizacion_maxima=1245.7500
\end{lstlisting}

\textit{Compilación y ejecución:}

\begin{lstlisting}[style=shellstyle]
# Desde practico-04/ejercicio-02
mkdir -p out
javac -d out src/main/java/ar/edu/unrc/bd2/practico4/ej2/*.java

# Interactivo (pide cod_banco por consola)
java -cp "out:lib/*" \
    ar.edu.unrc.bd2.practico4.ej2.CotizacionMaximaApp \
    config/db.postgres.properties

# O con cod_banco como argumento
java -cp "out:lib/*" \
    ar.edu.unrc.bd2.practico4.ej2.CotizacionMaximaApp \
    config/db.postgres.properties 1

# Variante MySQL
java -cp "out:lib/*" \
    ar.edu.unrc.bd2.practico4.ej2.CotizacionMaximaApp \
    config/db.mysql.properties 1
\end{lstlisting}

% =====================================================
\section{Ejercicio 3: Listado de metadatos (PK, UQ, FK)}
% =====================================================

\bloqueConsigna

Realizar un programa Java que liste las tablas de la base de datos del ejercicio 1 c) de la práctica 1) de repaso de SQL (bases de datos cliente, automóviles, etc), debe listar las tablas, sus columnas, clave primaria, claves únicas y foráneas.

\textbf{Nota:} Resolver utilizando los motores que se utilizaron en la creación de las bases de datos.

\bloqueIdea

El programa no se ocupa de datos, sino de \emph{metadatos}: la estructura de las tablas. JDBC expone esta información en \texttt{DatabaseMetaData}, una interfaz que estandariza los \texttt{ResultSet} con columnas como \texttt{TABLE\_NAME}, \texttt{COLUMN\_NAME}, \texttt{PK\_NAME}, \texttt{FKCOLUMN\_NAME}, etc. Eso permite escribir un programa portable entre motores: lo único que cambia entre PostgreSQL, Oracle y MySQL es el filtro de \texttt{catalog}/\texttt{schemaPattern}, que viaja en el \texttt{.properties}.

\bloqueSolucion

\textit{Estructura:}

\begin{itemize}
    \item \texttt{MetadataExplorerApp}: \texttt{main} mínimo; abre conexión, obtiene \texttt{DatabaseMetaData}, delega en el reporter.
    \item \texttt{MetadataReporter}: cinco pasos: cargar tablas, columnas, PK, UQ, FK.
    \item \texttt{DbConfig}: filtros para \texttt{getTables} (catalog, schemaPattern, tableNamePattern, tableTypes).
\end{itemize}

\textit{Configuración PostgreSQL:}

\begin{lstlisting}[style=propsstyle]
# config/db.postgres.properties
db.driver=org.postgresql.Driver
db.url=jdbc:postgresql://127.0.0.1:5432/postgres
db.user=postgres
db.password=

db.catalog=
db.schemaPattern=practico1c
db.tableNamePattern=%
db.table.types=TABLE
\end{lstlisting}

\textit{Configuración Oracle:}

\begin{lstlisting}[style=propsstyle]
# config/db.oracle.properties
db.driver=oracle.jdbc.OracleDriver
db.url=jdbc:oracle:thin:@//localhost:1521/XE
db.user=PRACTICO1C
db.password=PRACTICO1C

# En Oracle, el "schema" es el usuario; debe ir en MAYUSCULAS.
db.catalog=
db.schemaPattern=PRACTICO1C
db.tableNamePattern=%
db.table.types=TABLE
\end{lstlisting}

\textit{Carga de tablas:}

\begin{lstlisting}
import java.sql.DatabaseMetaData;
import java.sql.ResultSet;
import java.sql.SQLException;
import java.util.ArrayList;
import java.util.Comparator;
import java.util.List;

private List<TableRef> loadTables(DatabaseMetaData metaData, DbConfig config)
        throws SQLException {
    List<TableRef> tables = new ArrayList<>();
    try (ResultSet rs = metaData.getTables(
            config.getCatalog(),
            config.getSchemaPattern(),
            config.getTableNamePattern(),
            config.getTableTypes())) {
        while (rs.next()) {
            tables.add(new TableRef(
                rs.getString("TABLE_CAT"),
                rs.getString("TABLE_SCHEM"),
                rs.getString("TABLE_NAME")
            ));
        }
    }
    tables.sort(Comparator
        .comparing((TableRef t) -> nullSafe(t.schema))
        .thenComparing(t -> nullSafe(t.name)));
    return tables;
}
\end{lstlisting}

\textit{Columnas:}

\begin{lstlisting}
private void printColumns(DatabaseMetaData metaData, TableRef table) throws SQLException {
    try (ResultSet rs = metaData.getColumns(table.catalog, table.schema, table.name, "%")) {
        while (rs.next()) {
            String columnName = rs.getString("COLUMN_NAME");
            String typeName = rs.getString("TYPE_NAME");
            int dataType = rs.getInt("DATA_TYPE");
            int size = rs.getInt("COLUMN_SIZE");
            int decimalDigits = rs.getInt("DECIMAL_DIGITS");
            String nullable = rs.getString("IS_NULLABLE");

            String typeDescription = buildTypeDescription(typeName, dataType, size, decimalDigits);
            String nullDescription = "YES".equalsIgnoreCase(nullable) ? "NULL" : "NOT NULL";

            System.out.println("- " + columnName + " : " + typeDescription
                + " (" + nullDescription + ")");
        }
    }
}
\end{lstlisting}

\textit{Clave primaria (con \texttt{KEY\_SEQ} para PK compuestas):}

\begin{lstlisting}
import java.util.Map;
import java.util.TreeMap;

private void printPrimaryKey(DatabaseMetaData metaData, TableRef table) throws SQLException {
    Map<Short, String> pkColumns = new TreeMap<>();
    String pkName = null;

    try (ResultSet rs = metaData.getPrimaryKeys(table.catalog, table.schema, table.name)) {
        while (rs.next()) {
            pkName = rs.getString("PK_NAME");
            short keySeq = rs.getShort("KEY_SEQ");
            pkColumns.put(keySeq, rs.getString("COLUMN_NAME"));
        }
    }

    if (pkColumns.isEmpty()) {
        System.out.println("PK:\n- (sin clave primaria)");
        return;
    }
    System.out.println("PK:\n- " + (pkName != null ? pkName : "(sin nombre)")
        + " " + new ArrayList<>(pkColumns.values()));
}
\end{lstlisting}

\textit{Claves únicas (no hay \texttt{getUniqueKeys}: se usa \texttt{getIndexInfo} y se filtran las que coinciden con la PK):}

\begin{lstlisting}
try (ResultSet rs = metaData.getIndexInfo(table.catalog, table.schema, table.name,
        /* unique = */ true, /* approximate = */ false)) {
    while (rs.next()) {
        short type = rs.getShort("TYPE");
        if (type == DatabaseMetaData.tableIndexStatistic) continue;

        String indexName = rs.getString("INDEX_NAME");
        String columnName = rs.getString("COLUMN_NAME");
        short ordinalPosition = rs.getShort("ORDINAL_POSITION");
        uniqueIndexes.computeIfAbsent(indexName, k -> new TreeMap<>())
            .put(ordinalPosition, columnName);
    }
}
// Despues se descartan los indices unicos cuyas columnas coinciden con la PK.
\end{lstlisting}

\textit{Claves foráneas:}

\begin{lstlisting}
try (ResultSet rs = metaData.getImportedKeys(table.catalog, table.schema, table.name)) {
    while (rs.next()) {
        String fkName = rs.getString("FK_NAME");
        short keySeq  = rs.getShort("KEY_SEQ");
        FkPart part = new FkPart(
            rs.getString("FKCOLUMN_NAME"),
            rs.getString("PKTABLE_SCHEM"),
            rs.getString("PKTABLE_NAME"),
            rs.getString("PKCOLUMN_NAME")
        );
        groupedFks.computeIfAbsent(fkName, k -> new TreeMap<>()).put(keySeq, part);
    }
}
\end{lstlisting}

\textit{Punto de entrada:}

\begin{lstlisting}
import java.io.IOException;
import java.sql.Connection;
import java.sql.DatabaseMetaData;
import java.sql.SQLException;

public static void main(String[] args) {
    String configPath = args.length > 0 ? args[0] : "config/db.postgres.properties";
    DbConfig config;
    try {
        config = DbConfig.fromFile(configPath);
    } catch (IOException | IllegalArgumentException e) {
        System.err.println("Error cargando config: " + e.getMessage());
        System.exit(1); return;
    }

    try (Connection connection = config.openConnection()) {
        DatabaseMetaData metaData = connection.getMetaData();
        System.out.println("Conexion OK -> "
            + metaData.getDatabaseProductName() + " "
            + metaData.getDatabaseProductVersion());

        new MetadataReporter().printMetadata(metaData, config);

    } catch (ClassNotFoundException e) {
        System.err.println("Driver no encontrado: " + e.getMessage());
    } catch (SQLException e) {
        System.err.println("Error SQL: " + e.getMessage());
    }
}
\end{lstlisting}

\bloquePorque

\begin{itemize}
    \item \texttt{getTables(catalog, schemaPattern, tableNamePattern, types)} devuelve un \texttt{ResultSet} con las tablas filtradas. Las columnas (\texttt{TABLE\_CAT}, \texttt{TABLE\_SCHEM}, \texttt{TABLE\_NAME}, \texttt{TABLE\_TYPE}) están definidas por el estándar JDBC, así que el código no depende del motor.
    \item \texttt{getColumns} devuelve cada columna con \texttt{DATA\_TYPE} (código de \texttt{java.sql.Types}), \texttt{TYPE\_NAME} (nombre nativo, ej. \texttt{VARCHAR2}, \texttt{NUMBER}), \texttt{COLUMN\_SIZE}, \texttt{DECIMAL\_DIGITS} y \texttt{IS\_NULLABLE} (``YES''/``NO'').
    \item \texttt{getPrimaryKeys} puede devolver varias filas si la PK es compuesta. \texttt{KEY\_SEQ} indica el orden dentro de la PK; por eso se acumula en un \texttt{TreeMap<Short, String>} para ordenar automáticamente.
    \item JDBC no tiene \texttt{getUniqueKeys}: se piden los índices únicos con \texttt{getIndexInfo(..., unique=true, ...)} y se descartan los que coinciden exactamente con la PK (que también genera un índice único).
    \item \texttt{getImportedKeys} devuelve las FK \emph{salientes} de la tabla (las que esta tabla declara hacia otras); \texttt{getExportedKeys} sería el camino inverso. Una FK también puede ser compuesta, así que \texttt{KEY\_SEQ} se usa de nuevo.
    \item Los filtros varían por motor: en PostgreSQL, \texttt{schemaPattern = 'practico1c'}; en Oracle no hay esquema separado del usuario, así que va \texttt{'PRACTICO1C'} en mayúsculas; en MySQL no hay esquema sino \emph{catalog}, así que \texttt{catalog = 'practico1c\_mysql'} y \texttt{schemaPattern = null}.
\end{itemize}

\bloqueMySQL

El código Java de \texttt{DatabaseMetaData} es exactamente el mismo: lo único que cambia es el \texttt{.properties}. La diferencia conceptual a tener clara es que MySQL usa el slot \texttt{catalog} de JDBC para el nombre de la base, y \texttt{schemaPattern} queda en blanco (o \texttt{null}):

\begin{lstlisting}[style=propsstyle]
# config/db.mysql.properties
db.driver=com.mysql.cj.jdbc.Driver
db.url=jdbc:mysql://localhost:3306/practico1c_mysql?useSSL=false&serverTimezone=UTC&nullCatalogMeansCurrent=false
db.user=root
db.password=root

# En MySQL, el "catalog" es la base; el "schema" no aplica.
db.catalog=practico1c_mysql
db.schemaPattern=
db.tableNamePattern=%
db.table.types=TABLE
\end{lstlisting}

\textit{Particularidades de MySQL relevantes para metadatos:}

\begin{itemize}
    \item \texttt{nullCatalogMeansCurrent=false} en la URL evita el comportamiento heredado donde \texttt{catalog=null} matcheaba la base actual: con \texttt{false}, MySQL respeta el filtro tal como lo manda el cliente.
    \item En \texttt{getColumns}, el \texttt{TYPE\_NAME} aparece como \texttt{INT}, \texttt{VARCHAR}, \texttt{DECIMAL}, \texttt{DATETIME}, etc.; nada de \texttt{VARCHAR2} ni \texttt{NUMBER} como en Oracle, ni \texttt{INT4} como devuelve PostgreSQL.
    \item Para identificadores con caracteres reservados, MySQL usa backticks \texttt{`...`}; PostgreSQL y Oracle usan comillas dobles \texttt{"..."}. El programa no concatena identificadores ``a mano'', así que esto no afecta al código del Ejercicio 3.
    \item Para columnas \texttt{AUTO\_INCREMENT}, MySQL las marca con \texttt{IS\_AUTOINCREMENT = YES} en el \texttt{ResultSet} de \texttt{getColumns} (PostgreSQL hace lo mismo con \texttt{IDENTITY}; Oracle 10/11 no, porque no tiene \texttt{IDENTITY} hasta 12c).
\end{itemize}

\bloquePrueba

\begin{lstlisting}[style=shellstyle]
# Desde practico-04/ejercicio-03
mkdir -p out
javac -d out src/main/java/ar/edu/unrc/bd2/practico4/ej3/*.java

# PostgreSQL
java -cp "out:lib/*" \
    ar.edu.unrc.bd2.practico4.ej3.MetadataExplorerApp \
    config/db.postgres.properties

# Oracle
java -cp "out:lib/*" \
    ar.edu.unrc.bd2.practico4.ej3.MetadataExplorerApp \
    config/db.oracle.properties

# MySQL
java -cp "out:lib/*" \
    ar.edu.unrc.bd2.practico4.ej3.MetadataExplorerApp \
    config/db.mysql.properties
\end{lstlisting}

\textit{Salida esperada (PostgreSQL):}

\begin{lstlisting}[style=shellstyle]
Conexion OK -> PostgreSQL 14.10
============================================================
Tabla: practico1c.cliente
============================================================
Columnas:
- dni : INT4 (NOT NULL)
- nombre : VARCHAR(60) (NOT NULL)
- apellido : VARCHAR(60) (NOT NULL)
- direccion : VARCHAR(120) (NULL)
- tarifa : INT4 (NOT NULL)
PK:
- pk_cliente [dni]
Claves unicas:
- (sin claves unicas adicionales)
Claves foraneas:
- (sin claves foraneas)

============================================================
Tabla: practico1c.accidente
============================================================
...
Claves foraneas:
- fk_acc_cliente   [dni]        -> practico1c.cliente   [dni]
- fk_acc_taller    [nro_taller] -> practico1c.taller    [nro_taller]
- fk_acc_automovil [patente]    -> practico1c.automovil [patente]
\end{lstlisting}

% =====================================================
\section{Decisiones transversales}
% =====================================================

\begin{itemize}
    \item \textbf{Sin frameworks ORM.} Se usa únicamente el paquete \texttt{java.sql} (\texttt{Connection}, \texttt{PreparedStatement}, \texttt{CallableStatement}, \texttt{ResultSet}, \texttt{DatabaseMetaData}). El enunciado pide JDBC: nada de Hibernate, JPA, Spring Data.
    \item \textbf{Recursos cerrados con \texttt{try-with-resources}.} El compilador genera el \texttt{finally} con \texttt{close()}, incluso si el cuerpo lanza excepción. Es el patrón seguro desde Java 7.
    \item \textbf{Configuración externa.} \texttt{driver}, \texttt{url}, \texttt{user}, \texttt{password} y los nombres de tabla/función viven en \texttt{.properties}. El binario es independiente del entorno.
    \item \textbf{Independencia de motor donde tiene sentido.} El Ejercicio 1 cambia entre MySQL y PostgreSQL apenas cambiando el \texttt{.properties}; el Ejercicio 3 lo hace entre PostgreSQL, Oracle y MySQL.
    \item \textbf{Sentencias parametrizadas.} Todos los \texttt{INSERT}/\texttt{SELECT}/\texttt{CALL} usan \texttt{?} y \texttt{PreparedStatement}/\texttt{CallableStatement}. Cero concatenación, cero ventana para SQL Injection.
    \item \textbf{Errores diferenciados.} Cada \texttt{main} distingue \texttt{ClassNotFoundException} (driver ausente), \texttt{IOException} (problema con el \texttt{.properties}) y \texttt{SQLException} (problema con la base o la sentencia).
\end{itemize}

\section*{Conclusión}
\addcontentsline{toc}{section}{Conclusión}

Los tres ejercicios cubren los tres patrones canónicos de uso de JDBC: \emph{ABM con \texttt{PreparedStatement}}, \emph{invocación de stored procedure con \texttt{CallableStatement}} y \emph{exploración de metadatos con \texttt{DatabaseMetaData}}. La estructura repetida (\texttt{App} + \texttt{Service}/\texttt{Dao} + \texttt{DbConfig} + \texttt{.properties}) deja claro cómo cambia la API según el caso de uso, y los snippets ``En MySQL 8+'' agregan la portabilidad mínima necesaria para correr cada ejercicio también sobre MySQL: cambiar el driver, el formato de URL (\texttt{useSSL}, \texttt{serverTimezone}), y entender que MySQL usa \emph{catalog} donde PostgreSQL y Oracle usan \emph{schema}.

\end{document}
